\chapter{Feature engineering}
\label{ch:feature-engineering}

\begin{quote}
\emph{``Coming up with features is difficult, time-consuming, requires expert knowledge. Applied machine learning is basically feature engineering.''}
--- Andrew Ng
\end{quote}

% ============================================================
\section{What is feature engineering?}

Feature engineering is the process of creating new input variables from existing data to improve model performance. While encoding and scaling (Chapter~5) transform existing features, feature engineering \emph{creates} new ones that capture patterns the raw data does not explicitly represent.

\begin{datatip}
A well-engineered feature can be worth more than a complex model. A simple linear regression with the right features often outperforms a deep neural network with raw features.
\end{datatip}

% ============================================================
\section{Polynomial and interaction features}

\begin{minted}{python}
import pandas as pd
import numpy as np
from sklearn.preprocessing import PolynomialFeatures

# Melbourne Housing dataset
df = pd.read_csv("melb_data.csv")
df = df.dropna(subset=["Price", "Rooms", "Distance", "Landsize"])

# Polynomial features: capture non-linear relationships
poly = PolynomialFeatures(degree=2, interaction_only=False,
                          include_bias=False)
features = df[["Rooms", "Distance"]]
poly_features = poly.fit_transform(features)
poly_names = poly.get_feature_names_out(["Rooms", "Distance"])

poly_df = pd.DataFrame(poly_features, columns=poly_names)
print(poly_df.head())
print(f"Original features: {features.shape[1]}")
print(f"Polynomial features: {poly_df.shape[1]}")
\end{minted}

\begin{minted}{python}
# Interaction only: no squared terms, just cross-products
poly_inter = PolynomialFeatures(degree=2, interaction_only=True,
                                 include_bias=False)
inter_features = poly_inter.fit_transform(features)
inter_names = poly_inter.get_feature_names_out(["Rooms", "Distance"])
print(f"Interaction features: {inter_names}")
\end{minted}

\begin{warningbox}
Polynomial features grow combinatorially: $n$ features at degree $d$ produce $\binom{n+d}{d} - 1$ features. With 20 features at degree~3, you get 1770 features. Use this technique selectively, not on the entire feature set.
\end{warningbox}

% ============================================================
\section{Mathematical transformations}

\begin{minted}{python}
# Ratio features: often more informative than raw values
df["price_per_room"] = df["Price"] / df["Rooms"]
df["price_per_sqm"] = df["Price"] / df["Landsize"].replace(0, np.nan)
df["rooms_per_bathroom"] = df["Rooms"] / df["Bathroom"].replace(0, 1)

print(df[["Price", "Rooms", "price_per_room",
          "price_per_sqm"]].describe())
\end{minted}

\begin{minted}{python}
# Log transform for skewed features
df["log_price"] = np.log1p(df["Price"])
df["log_landsize"] = np.log1p(df["Landsize"])

# Square root: gentler than log
df["sqrt_distance"] = np.sqrt(df["Distance"])

# Compare correlation with target
for col in ["Distance", "sqrt_distance", "log_price"]:
    if col in df.columns and df[col].notnull().sum() > 0:
        corr = df[col].corr(df["Price"])
        print(f"Corr({col}, Price) = {corr:.3f}")
\end{minted}

\begin{preprocesstip}
Ratio features encode \emph{domain knowledge} in a single number. ``Price per square meter'' is a standard real estate metric that captures value density better than price or area alone.
\end{preprocesstip}

% ============================================================
\section{Date and time features}

\begin{minted}{python}
# Melbourne Housing: extract date components
df["Date"] = pd.to_datetime(df["Date"], dayfirst=True)

df["sale_year"] = df["Date"].dt.year
df["sale_month"] = df["Date"].dt.month
df["sale_dayofweek"] = df["Date"].dt.dayofweek  # 0=Monday
df["sale_quarter"] = df["Date"].dt.quarter
df["sale_is_weekend"] = df["Date"].dt.dayofweek.isin([5, 6]).astype(int)

print(df[["Date", "sale_year", "sale_month",
          "sale_dayofweek", "sale_quarter"]].head())
\end{minted}

\begin{minted}{python}
# Cyclical encoding for month (so December and January are close)
df["month_sin"] = np.sin(2 * np.pi * df["sale_month"] / 12)
df["month_cos"] = np.cos(2 * np.pi * df["sale_month"] / 12)

# Days since a reference date
reference = pd.Timestamp("2016-01-01")
df["days_since_ref"] = (df["Date"] - reference).dt.days

print(df[["sale_month", "month_sin", "month_cos"]].head(12))
\end{minted}

% ============================================================
\section{Aggregation features}

\begin{minted}{python}
# Suburb-level statistics as features
suburb_stats = df.groupby("Suburb")["Price"].agg(
    suburb_mean_price="mean",
    suburb_median_price="median",
    suburb_price_std="std",
    suburb_n_sales="count"
).reset_index()

df = df.merge(suburb_stats, on="Suburb", how="left")

# How does this property compare to its suburb average?
df["price_vs_suburb"] = df["Price"] / df["suburb_mean_price"]

print(df[["Suburb", "Price", "suburb_mean_price",
          "price_vs_suburb"]].head(10))
\end{minted}

\begin{minted}{python}
# Gapminder: country-level features over time
gap = pd.read_csv("gapminder.csv")

# Year-over-year growth rate
gap = gap.sort_values(["country", "year"])
gap["gdp_growth"] = gap.groupby("country")["gdpPercap"].pct_change()
gap["pop_growth"] = gap.groupby("country")["pop"].pct_change()

# Rolling average (smoothed trend)
gap["lifeExp_rolling3"] = (gap.groupby("country")["lifeExp"]
                               .transform(lambda x: x.rolling(3).mean()))

print(gap[gap["country"] == "Benin"][
    ["year", "lifeExp", "lifeExp_rolling3", "gdp_growth"]].tail(6))
\end{minted}

% ============================================================
\section{Text-derived features}

\begin{minted}{python}
# Extract features from text columns without full NLP
# Using Melbourne suburb names and addresses

# Length of address
df["address_length"] = df["Address"].str.len()
df["address_word_count"] = df["Address"].str.split().str.len()

# Contains specific keywords
df["is_street"] = df["Address"].str.contains("St$", regex=True).astype(int)
df["is_road"] = df["Address"].str.contains("Rd$", regex=True).astype(int)
df["is_avenue"] = df["Address"].str.contains("Av", regex=False).astype(int)

print(df[["Address", "address_length", "is_street",
          "is_road"]].head(10))
\end{minted}

% ============================================================
\section{Domain knowledge features}

\begin{minted}{python}
# Real estate domain knowledge
df["total_rooms"] = df["Rooms"] + df["Bathroom"] + df["Car"]
df["has_multiple_bathrooms"] = (df["Bathroom"] > 1).astype(int)
df["is_new_build"] = (df["YearBuilt"] > 2010).astype(int)
df["building_age"] = df["sale_year"] - df["YearBuilt"]
df["is_close_to_cbd"] = (df["Distance"] < 10).astype(int)

# Property type indicators
df["is_house"] = (df["Type"] == "h").astype(int)
df["is_unit"] = (df["Type"] == "u").astype(int)

# Interaction: large house close to CBD
df["large_close"] = df["is_house"] * df["is_close_to_cbd"] * df["Rooms"]

print(df[["Suburb", "Type", "Rooms", "Distance",
          "building_age", "large_close"]].head(10))
\end{minted}

\begin{datatip}
The best features come from understanding the domain. Spend time with domain experts. Ask: ``What factors do professionals in this field consider when making decisions?'' Encode those factors as features.
\end{datatip}

% ============================================================
\section{Feature selection after engineering}

\begin{minted}{python}
from sklearn.feature_selection import mutual_info_regression

# Select numerical features
feature_cols = ["Rooms", "Distance", "Landsize", "building_age",
                "price_per_room", "suburb_mean_price", "large_close",
                "days_since_ref", "Bathroom", "Car"]
target = "Price"

# Drop rows with NaN for this analysis
subset = df[feature_cols + [target]].dropna()

mi = mutual_info_regression(subset[feature_cols], subset[target],
                            random_state=42)
mi_series = pd.Series(mi, index=feature_cols).sort_values(ascending=False)

print("Mutual information with Price:")
print(mi_series)
\end{minted}

\begin{minted}{python}
# Correlation matrix for engineered features
import seaborn as sns
import matplotlib.pyplot as plt

corr = subset[feature_cols].corr()
fig, ax = plt.subplots(figsize=(10, 8))
sns.heatmap(corr, annot=True, fmt=".2f", cmap="coolwarm",
            center=0, ax=ax)
plt.title("Feature correlation matrix")
plt.tight_layout()
plt.savefig("feature_correlation.png", dpi=150)
plt.show()
\end{minted}

% ============================================================
\section{Exercises}

\begin{exercise}
\begin{enumerate}
  \item For the Melbourne Housing dataset, create at least 5 new features using domain knowledge. Rank them by mutual information with Price.
  \item Create cyclical encoding for the sale month. Verify that December and January are close in the encoded space by computing their Euclidean distance.
  \item Using the Gapminder dataset, create year-over-year growth rate features for GDP per capita and life expectancy. Which country had the highest GDP growth in a single period?
  \item Build polynomial features (degree~2) for \texttt{Rooms}, \texttt{Distance}, and \texttt{Landsize}. Train a linear regression with and without polynomial features and compare $R^2$.
  \item Create a feature engineering function \texttt{engineer\_melb(df)} that takes the raw Melbourne dataset and returns a DataFrame with at least 10 new features, ready for modeling.
\end{enumerate}
\end{exercise}

% ============================================================
\section{Chapter summary}

\begin{itemize}
  \item Feature engineering creates new variables that capture patterns invisible in raw data.
  \item Polynomial and interaction features capture non-linear and cross-feature relationships.
  \item Ratio features (price per room, price per square meter) encode domain knowledge compactly.
  \item Date features (year, month, day of week, cyclical encoding) unlock temporal patterns.
  \item Aggregation features (group means, rolling averages) add contextual information.
  \item Domain knowledge is the most valuable source of feature ideas --- consult experts.
  \item After engineering, use mutual information or correlation to select the most predictive features.
\end{itemize}
